\documentclass[11pt]{article}
\usepackage{fontspec,amsmath,amssymb,amsthm,hyperref}
\usepackage[margin=1in]{geometry}
\setlength{\emergencystretch}{3em}
\title{Gaussian fourth moments control the actual Taylor error}
\author{LangevinLean mathematical exposition}
\date{}
\begin{document}\maketitle
\noindent Source module: \texttt{KineticLangevin/LpMomentBounds.lean}. Compile this proof document with XeLaTeX or LuaLaTeX. The accompanying HTML pairs the mathematical argument with the exact Lean statements and records the scope of the formalization.


The quadratic Hessian remainder requires a fourth displacement moment, even though the stochastic solution itself lives in L\ensuremath{{}^{2}}. These results connect actual L4 seminorms, exact Gaussian fourth moments, and the existing L\ensuremath{{}^{2}} kinematic equation.


\section{Turn a pointwise quadratic remainder into an L\ensuremath{{}^{2}} bound}


Suppose actual random vectors D and X satisfy ‖D(ω)‖\ensuremath{\le}a‖X(ω)‖\ensuremath{{}^{2}}, with a≥0 and ‖X‖₄\ensuremath{\le}C. Monotonicity of the extended L\ensuremath{{}^{2}} norm gives ‖D‖\ensuremath{{}_{2}}\ensuremath{\le}a‖‖X‖\ensuremath{{}^{2}}‖\ensuremath{{}_{2}}. The power rule identifies the last factor with ‖X‖₄\ensuremath{{}^{2}}, yielding aC\ensuremath{{}^{2}}. Lean works first with extended nonnegative norms, where the power and monotonicity rules hold without hidden finiteness assumptions, and then converts the finite upper bound back to the actual L\ensuremath{{}^{2}} norm.



\noindent Lean declarations: \texttt{lp\_norm\_le\_quadratic}.

\section{Compute the actual Gaussian L4 radius}


A standard Gaussian vector in p dimensions has fourth norm moment p\ensuremath{{}^{2}}+2p. Fernique integrability ensures its fourth seminorm is finite. The integral formula for the L4 seminorm therefore yields (p\ensuremath{{}^{2}}+2p)\textasciicircum{}(1/4) on any realization of the Gaussian law. This proof transfers an actual integral identity through HasLaw and keeps the 2p correction, including the one-dimensional case.



\noindent Lean declarations: \texttt{gaussian\_eLpNorm\_four}.

\section{Apply it to the existing kinetic displacement}


Assume X(t)−X(0)=∫₀ᵗV(s)ds in the actual L\ensuremath{{}^{2}} space and every V(s) has the standard Gaussian marginal. Each velocity lies in the same L4 ball of radius (p\ensuremath{{}^{2}}+2p)\textasciicircum{}(1/4). The closed-convex averaging theorem for the existing L\ensuremath{{}^{2}} integral gives ‖X(t)−X(0)‖₄\ensuremath{\le}t(p\ensuremath{{}^{2}}+2p)\textasciicircum{}(1/4). No independence across times is used or required. This supplies exactly the fourth moment needed by the Hessian Taylor remainder in Proposition 2.



\noindent Lean declarations: \texttt{gaussian\_kinematic\_four\_bound}.

\end{document}
